\section{Handling Imbalanced Pattern Discovery}
\label{sec:2_5}

Fraud detection datasets exhibit extreme class imbalance, with fraudulent instances typically comprising 0.1-5\% of observations. This skew fundamentally affects pattern discovery: minority class patterns may be overlooked when learning is dominated by majority class examples.

\subsection{Class Imbalance in Fraud Data}

Class imbalance creates two interrelated challenges for pattern mining.

\textbf{Learning bias} causes models to optimize for majority class patterns. Standard loss functions weight all samples equally; with 1\% fraud rate, a model predicting ``legitimate'' for all cases achieves 99\% accuracy while discovering zero fraud patterns. Neural network training, including GNNs, converges toward trivial solutions under extreme imbalance \cite{buda2018systematic}. Message passing from majority-class neighbors further reinforces legitimate patterns while suppressing fraud signal \cite{liu2021heterogeneous}.

\textbf{Pattern sparsity} means minority class patterns may lack sufficient examples for robust discovery. Fraud rings vary in structure; learning generalizable patterns requires diverse fraud examples. \citeauthor{he2009learning} \cite{he2009learning} surveyed approaches to imbalanced learning, categorizing solutions into data-level (resampling), algorithm-level (cost-sensitive), and hybrid methods.

\enlargethispage{2\baselineskip}
\subsection{Sampling Strategies}

\textbf{Oversampling} methods increase minority class representation through synthetic example generation. SMOTE \cite{chawla2002smote} creates synthetic samples via interpolation between minority instances and their k-nearest neighbors. Borderline-SMOTE \cite{han2005borderline} focuses generation on decision boundary examples. ADASYN \cite{he2008adasyn} weights generation by local difficulty, creating more samples in regions with higher misclassification risk.

\textbf{Critical methodological requirement:} Resampling must occur exclusively within training folds during cross-validation. Applying SMOTE before train-test splitting causes synthetic samples to encode information from test instances, inflating performance estimates. \citeauthor{demircioglu2024measuring} \cite{demircioglu2024measuring} quantified this leakage at +0.34 AUC bias sufficient to invalidate experimental conclusions.

\textbf{Cost-sensitive learning} adjusts loss functions rather than data distributions. XGBoost's \texttt{scale\_pos\_weight} parameter multiplies positive class loss, theoretically balancing class contributions. Focal Loss \cite{lin2017focal} down-weights well-classified examples regardless of class, focusing learning on difficult cases where fraud patterns are ambiguous.

\enlargethispage{\baselineskip}
\subsection{Evaluation for Imbalanced Patterns}

Metric selection critically affects pattern validation conclusions.

\textbf{AUC-ROC} measures discrimination across thresholds but incorporates true negative rate, which is inflated when majority class dominates. A model can achieve AUC-ROC = 0.95 while detecting minimal fraud at practical thresholds.

\textbf{AUC-PR} (Area Under Precision-Recall curve) focuses on positive class performance. \citeauthor{davis2006relationship} \cite{davis2006relationship} proved that PR curves provide more informative visualization for skewed datasets. \citeauthor{saito2015precision} \cite{saito2015precision} provided empirical evidence that AUC-PR reveals performance differences invisible in ROC space.

\textbf{Precision@K} directly addresses operational constraints what fraction of top-K flagged items are genuinely fraudulent? This metric maps to analyst capacity and investigation efficiency.

This thesis adopts AUC-PR as primary metric, with Precision@K for operational interpretation.
